\documentclass[10pt,a4paper]{amsart}
\usepackage[margin=19mm]{geometry}
\usepackage{amsmath,amssymb,mathtools}
\usepackage[scr=rsfs,cal=euler]{mathalfa}
\usepackage{xfrac}
\usepackage[unicode,hidelinks,bookmarks=false]{hyperref}
\newcommand{\ie}{i.e.\ }
\newcommand{\cf}{cf.\ }
\newcommand{\resp}{respectively\ }
\newcommand{\Subsection}{\subsection}
\providecommand{\nobreakdash}{\nobreak\hbox{-}}
\setlength{\emergencystretch}{2em}
\multlinegap0pt
\usepackage{bm}
\usepackage{bbm}
\usepackage{dsfont}
\usepackage{xspace}
\usepackage{cancel}
\usepackage{mathtools}
\usepackage{amsthm}
\makeatletter
\@ifundefined{lemma}{\newtheorem{lemma}{Lemma}}{}
\@ifundefined{proposition}{\newtheorem{proposition}[lemma]{Proposition}}{}
\makeatother
\newcommand{\e}{\varepsilon} 
\title[Higher-order non-local gradient theory of phase-transitions]{Higher-order non-local gradient theory of phase-transitions}
\author{Margherita Solci}
\date{Source: arXiv:2411.01586v2 (CC BY 4.0)}
\begin{document}
\maketitle

\noindent\emph{Source and licence.} Higher-order non-local gradient theory of phase-transitions. Version arXiv:2411.01586v2 (CC BY 4.0), distributed under the Creative Commons Attribution 4.0 International licence, as recorded in the arXiv licence metadata for this exact version and shown on the abstract page https://arxiv.org/abs/2411.01586v2. Full rights notice: licenses/arxiv-2411.01586-CC-BY-4.0.txt.\par\smallskip

\noindent\textit{Editorial scope.} Counted unit: the Proposition whose source label is \texttt{None} in the article (source0.tex lines 1099--1101, source label \texttt{None}), delivered together with the article's own complete original proof of it (source0.tex lines 1103--1222, one \texttt{proof} environment of 120 lines). No mathematics is written, repaired or re-derived: statement and proof are the article's own text line for line. Required context retained uncounted: none. Every definition, display and label of those lines is retained unchanged. Omitted from this package and not redistributed: 1--1098, 1102--1102, 1223--1312. Nothing inside a retained excerpt is omitted or altered: every retained body is delivered exactly as the article's lines are written, so no omission receipt of any kind accompanies this package. Citations: the retained excerpts carry no citation command at all, so no bibliography entry of the article is transcribed and no reference is redistributed. Reference adaptation: none was needed and none was made; every reference inside the delivered unit resolves inside it, so no unresolved reference or citation is produced. Printed slips kept literal: no printed word, symbol, hypothesis or number of the counted statement or of its proof was altered. Layout and class: the article's own class is \texttt{article} with packages including amsmath, amsthm, amssymb, color, xcolor; the wrapper is the lane's pinned \texttt{amsart} editorial wrapper at 10pt on A4, which restates the theorem-like environments the retained text needs and adds the article's own macro definitions that the retained text needs (\texttt{\textbackslash e}), with extra wrapper packages (bm, bbm, dsfont, xspace, cancel, mathtools). The delivered numbering is the unit's own. Assets: no retained span contains a figure environment, a graphics-inclusion command, a PDF-page inclusion, a table or a TikZ picture; no image file is copied into this package, the package declares no asset dependency and no image file is redistributed with it. Licence: version arXiv:2411.01586v2 of this article is distributed under the Creative Commons Attribution 4.0 International licence, as recorded in the arXiv licence metadata for this exact version and shown on the abstract page https://arxiv.org/abs/2411.01586v2. A full rights notice is delivered at licenses/arxiv-2411.01586-CC-BY-4.0.txt. Characters: every retained excerpt is pure ASCII, so the delivered engine, XeLaTeX with its default Latin Modern text font, reports no missing character.
\begin{proposition}
If $u\colon(0,1)\to\{-1,1\}$ has a finite discontinuity set $S(u)$, then there exists a family $u_\e$ such that $u_\e\to u$ in $L^1(0,1)$ and such that $F_\e(u_\e)$ converges to $\widetilde m_{k+s}\#S(u)$.
\end{proposition}

\begin{proof}
Up to a diagonal argument, and taking the lower bound into account, it suffices to prove that for each $\sigma>0$ there exists a sequence $u_\e$ such that $u_\e\to u$ in $L^1(0,1)$ and 
\begin{equation}
\limsup_{\e\to 0}F_\e(u_\e)\le\widetilde m_{k+s}\#S(u)+\sigma.
\end{equation}

\smallskip
Let $\eta>0$ be fixed and let $T>0$ and $v\in H^{k+s}_{\rm loc} (\mathbb R)$ be such that 
$$
\int_{-\infty}^{+\infty} W(v)dx+[v^{(k)}]_{s}(\mathbb R)\le \widetilde m_{k+s}+\eta
$$
and $v(t)=1$ for $t\ge T$, $v(t)=-1$ for $t\le -T$.

If  $N=\#S(u)$ and $S(u)=\{t_1,\ldots, t_N\}$ with $t_i<t_{i+1}$, we define the sets
$$
S_+=\{i\in\{1,\ldots,N\}: u(t_i^+)=1\},\quad S_-=\{i\in\{1,\ldots,N\}: u(t_i^-)=1\}.
$$
Note that one of these two sets is just the set of even $i$ and the other the set of odd $i\in\{1,\ldots,N\}$.
For $\e$ small enough we define $u_\e\colon(0,1)\to \mathbb R$ as 
$$
u_\e(t)=\begin{cases}
v(\frac{t-t_i}\e) &\hbox{ if }|t-t_i|\le\e T, i\in S_+\\
v(\frac{t_i-t}\e) &\hbox{ if }|t-t_i|\le\e T, i\in S_-\\
u(t) &\hbox{ otherwise.}
\end{cases}
$$
The functions thus defined belong to $H^{k+s}(0,1)$. We now estimate $F_\e(u_\e)$.

If we set $x_0=0$, $x_i=\frac{t_i+t_{i+1}}2$ for $i\in\{1,\ldots, N-1\}$ and $x_N=1$, 
then we can write
\begin{eqnarray*}
&&F_\e(u_\e)=\sum_{i=1}^N\frac1\e\int_{t_i-\e T}^{t_i+\e T}W(u_\e)dx 
\\&&\hskip1.5cm+\e^{2(k+s)-1}
\sum_{i=1}^N\sum_{j=1}^N\int_{x_{i-1}}^{x_i}\int_{x_{j-1}}^{x_j}\frac{|u^{(k)}_\e(x)-
u^{(k)}_\e(y)|^2}{|x-y|^{1+2s}}dx\,dy.
\end{eqnarray*}
Changing variables, we have 
$$
\frac1\e\int_{t_i-\e T}^{t_i+\e T}W(u_\e)dt= \int_{-T}^T W(v)dt
$$
and 
$$
\e^{2(k+s)-1}\int_{x_{i-1}}^{x_i}\int_{x_{i-1}}^{x_i}\frac{|u^{(k)}_\e(x)-
u^{(k)}_\e(y)|^2}{|x-y|^{1+2s}}dx\,dy\le [v^{(k)}]_{s}(\mathbb R)
$$
 for all $i\in\{1,\ldots, N\}$.
Hence, the claim follows once we prove that each term
\begin{equation}\label{Iij} I^\e_{ij}:=
\e^{2(k+s)-1}\int_{x_{i-1}}^{x_i}\int_{x_{j-1}}^{x_j}\frac{|u^{(k)}_\e(x)-
u^{(k)}_\e(y)|^2}{|x-y|^{1+2s}}dy\,dx
\end{equation}
in the double sum with $i\neq j$ is asymptotically negligible. After a change of variables we can write
$$
I^\e_{ij}=\int_{\frac{x_{i-1}-t_i}\e}^{\frac{x_i-t_i}\e}\int_{\frac{x_{j-1}-t_j}\e}^{\frac{x_j-t_j}\e} \frac{|v^{(k)}(x)-
v^{(k)}(\pm y)|^2}{\big|y-x+\frac{t_j-t_i}\e \big|^{1+2s}}dy\,dx,
$$
with the plus sign if $i-j$ is even, and with the minus sign otherwise. The term $t_i-t_j$ in the denominator shows that the double integrals with the largest values are those with consecutive indices, then it suffices to treat the case when $i=1$ and $j=2$.  It is not restrictive to suppose that $1\in S_+$. We then have to estimate 
$$
I^\e=\int_{\frac{x_0-t_1}\e}^{\frac{x_1-t_1}\e}\int_{\frac{x_1-t_2}\e}^{\frac{x_2-t_2}\e} \frac{|v^{(k)}(x)-
v^{(k)}(-y)|^2}{\big|y-x+\frac{t_2-t_1}\e \big|^{1+2s}}dy\,dx.
$$
Since $v^{(k)}(x)=v^{(k)}(-y)$ if $x\ge T$ and $y\le -T$ or  if $x\le -T$ and $y\ge T$ (note that if $k\ge 1$ this is true also if $x,y\ge T$ or $x,y\le -T$), it suffices to estimate the double integrals
$$
I^\e_1:=\int_{-T}^{T}\int_{\frac{x_1-t_2}\e}^{\frac{x_2-t_2}\e} \frac{|v^{(k)}(x)-
v^{(k)}(-y)|^2}{\big|y-x+\frac{t_2-t_1}\e \big|^{1+2s}}dy\,dx
$$
and 
$$
I^\e_2:=\int_{T}^{\frac{x_1-t_1}\e}\int_{-T}^{\frac{x_2-t_2}\e} \frac{|v^{(k)}(x)-
v^{(k)}(-y)|^2}{\big|y-x+\frac{t_2-t_1}\e \big|^{1+2s}}dy\,dx,
$$
and the analogous $I_3$ and $I_4$ with the role of $x$ and $y$ interchanged.

We first estimate
\begin{eqnarray}\label{Ie1}\nonumber
I^\e_1&\le& 2\int_{-T}^{T}\int_{\frac{x_1-t_2}\e}^{\frac{x_2-t_2}\e} \frac{|v^{(k)}(x)|^2+
|v^{(k)}(-y)|^2}{\big|y-x+\frac{t_2-t_1}\e \big|^{1+2s}}dy\,dx
\\ \nonumber
&= &2\int_{-T}^{T}|v^{(k)}(x)|^2\Big(\int_{\frac{x_1-t_2}\e}^{\frac{x_2-t_2}\e} +
\frac1{\big|y-x+\frac{t_2-t_1}\e \big|^{1+2s}}dy\Big)\,dx
\\
&&+2\int_{-T}^{T}\Big(\int_{\frac{x_1-t_2}\e}^{\frac{x_2-t_2}\e} \frac{
|v^{(k)}(-y)|^2}{\big|y-x+\frac{t_2-t_1}\e \big|^{1+2s}}dy\Big)\,dx.
\end{eqnarray}
For $|x|\le T$ we have
$$
\int_{\frac{x_1-t_2}\e}^{\frac{x_2-t_2}\e} \frac{
1}{\big|y-x+\frac{t_2-t_1}\e \big|^{1+2s}}dy
\le\frac{\frac{t_2-t_1}\e}{\big|\frac{x_1-t_1}\e-T \big|^{1+2s}} =O(\e^{2s})
$$
as $\e\to 0$, and hence also
\begin{eqnarray*}
&&\hskip-1.5cm\int_{\frac{x_1-t_2}\e}^{\frac{x_2-t_2}\e} \frac{
|v^{(k)}(-y)|^2}{\big|y-x+\frac{t_2-t_1}\e \big|^{1+2s}}dy\ 
\le\ \int_{\frac{x_1-t_2}\e}^{-T} \frac{
1}{\big|y-x+\frac{t_2-t_1}\e \big|^{1+2s}}dy
\\
&&+\int_{-T}^{T} \frac{
|v^{(k)}(-y)|^2}{\big|y-x+\frac{t_2-t_1}\e \big|^{1+2s}}dy
+
\int_{T}^{\frac{x_2-t_2}\e} \frac{1}{\big|y-x+\frac{t_2-t_1}\e \big|^{1+2s}}dy
\\
&\le&
\int_{\frac{x_1-t_2}\e}^{\frac{x_2-t_2}\e} \frac{
1}{\big|y-x+\frac{t_2-t_1}\e \big|^{1+2s}}dy
+\int_{-T}^{T} \frac{
|v^{(k)}(-y)|^2}{\big|y-x+\frac{t_2-t_1}\e \big|^{1+2s}}dy
\\
&\le& \frac{\frac{x_2-x_1}\e}{\big|\frac{x_1-t_1}\e-T \big|^{1+2s}}+\frac1{\big|\frac{t_2-t_1}\e-2T \big|^{1+2s}}\int_{-T}^{T} 
|v^{(k)}(y)|^2dy=o(1)
\end{eqnarray*}
as $\e\to 0$. Using these estimates in \eqref{Ie1} we obtain that $I^\e_1=o(1)$ as $\e\to 0$. Note that the second estimate slightly improves if $k\ge 1$ since $v^{(k)}(y)=0$ if $|y|\ge T$. An analogous argument shows that $I^\e_2=o(1)$ as $\e\to 0$,
so that $I^\e$ tends to $0$ as $\e\to 0$. Note that if $i-j$ is even, the argument above can be followed word for word since we have only used the estimate $|v^{(k)}(x)-v^{(k)}(-y)|^2\le 2(|v^{(k)}(x)|^2+|v^{(k)}(-y)|^2)$, and hence the claim is proved.

We finally note that if $s>\frac12$ then the estimate of $I^\e_{ij}$ defined in \eqref{Iij}  for $i\neq j$ can be achieved by a simpler direct computation, as we have
$$
I^\e_{ij}\le\e^{2s-1} \int_{x_{i-1}}^{x_i}\int_{x_{j-1}}^{x_j}\frac{4\|v^{(k)}\|_\infty^2}{|x-y|^{1+2s}}dy\,dx\le \e^{2s-1}\frac{8\|v^{(k)}\|_\infty^2}{s(2s-1){\delta^{2s-1}}},
$$
where $\delta=\min\{|x_k-x_\ell|: k\neq \ell\}$.
\end{proof} 

\end{document}
